\begin{afterword}
\summaryhead

The post begins with two observations: every Utopia is to some degree a place no one would want to live, as Orwell also said, and books on writing insist that stories need conflict, or, in Orson Scott Card's and Jack Bickham's plainer words, pain and disaster. The author once thought that a positive Singularity would be the end of all stories, and later wondered whether ``no stories'' was a sign of a world not worth living in. The post examines the argument that a world without pain has no stories worth reading and so is not worth living in. It grants that stories and lives are optimized by different criteria, but asks whether the jarring quality of stories where everything goes right would carry over to life.

It then argues that pain is far easier to produce, and far more intense, than pleasure, which may be why the great stories are tragedies. The author guesses that pain and hedonic adaptation could be engineered away, and judges David Pearce ``very probably right'' that abolishing suffering is feasible. Against complete abolition the author sets ``the path of courage'': remove the pain that destroys people, make people stronger, and keep lesser sorrows. The author admits that this preference may come from fear, and ends: ``I don't know.''

\responsehead

The post is an exploration. Most of its claims are hedged, and it ends undecided. Read as that, it handles its sources well. The quotation from Orwell is exact, and his essay supports the opening claim. The words given to Card and Bickham match the substance of their manuals, though I could not find either as an exact quotation. The loss-aversion figure is near the mean of later estimates. The puzzle of why we prize stories of suffering is an old one, Hume's in ``Of Tragedy,'' and the claim that pain is more powerful than pleasure has support in the psychology of ``bad is stronger than good.''

On stories, the post does its own work. The argument it sets out depends on its second premise, that a world not worth reading about is not worth living in, and the post rejects that premise itself. Its real question is narrower: whether the discomfort of stories where everything goes right would carry over to life. Here the manuals it quotes point toward its own first answer, that stories need striving rather than pain. Card chooses a character in pain because such a character ``wants things to change'' and acts; Bickham's ``disaster'' is a setback to a character's goal, not an earthquake. The list of great tragedies is the author's choice and illustrates rather than tests the claim; Aristotle, who also ranked the unhappy ending first, noted that the happy one was preferred ``because of the weakness of the spectators.''

On whether pain and adaptation could be removed, the author guesses three times, then calls Pearce ``very probably right,'' ``So far as I know or can guess.'' No evidence about brains is given. The two facts the post does offer need qualification. The separation of the circuits for pleasure and pain is less clear than ``distinct'' suggests. People born without pain suffer bitten tongues, fractures and, in one reported case, an early death: a larger cost than having to inspect themselves.

The ending states a preference for ``the path of courage'' and names two biases that might explain it. As a preference it needs no evidence, and the post does not claim more for it.

In short: an open question about whether a good world needs pain, with sources that hold up and an old puzzle about tragedy at its center, which ends, undecided, with a stated preference for ``the path of courage.''
\end{afterword}
